\documentclass[11pt,a4paper]{article}
\usepackage{fontspec,amsmath,amssymb,amsthm,microtype}
\setmainfont{lmroman10-regular.otf}[BoldFont=lmroman10-bold.otf,ItalicFont=lmroman10-italic.otf,BoldItalicFont=lmroman10-bolditalic.otf]
\usepackage[margin=25mm]{geometry}
\usepackage[unicode,hidelinks]{hyperref}
\hypersetup{pdfauthor={LangevinLean Project}}
\setlength{\emergencystretch}{3em}
\setlength{\parskip}{0.4em}
\setlength{\parindent}{0pt}
\allowdisplaybreaks[2]
\raggedbottom
\title{\texorpdfstring{Exact Gaussian averaging of the quadratic Taylor polynomial}{Exact Gaussian averaging of the quadratic Taylor polynomial}}
\author{LangevinLean Project}
\date{}
\begin{document}\maketitle
Draft mathematical exposition; these notes may contain mistakes.


All expectations here use the actual standard Gaussian measure in a finite-dimensional real Hilbert space. The quadratic trace identity identifies the diffusion coefficient in the Langevin generator.



For an actual standard Gaussian vector and any continuous bilinear form,


\[\mathbb E[\langle u,Z\rangle\langle v,Z\rangle]=\langle u,v\rangle,
\qquad \mathbb E[B(Z,Z)]=\sum_i B(e_i,e_i),\qquad Z\sim N(0,I).\]


Writing the quadratic Taylor polynomial as \(\mathcal T_2\phi(x;w)=\phi(x)+D\phi(x)[w]+\tfrac12D^2\phi(x)[w,w]\), its Gaussian expectation is exactly


\[\mathbb E\,\mathcal T_2\phi\bigl(x;hv+\sqrt{2h}Z\bigr)
=\phi(x)+h\bigl(D\phi(x)[v]+\Delta\phi(x)\bigr)
+\frac{h^2}{2}D^2\phi(x)[v,v].\]

Earlier results used here: \cite[Section~3.55]{LangevinLean2026ULASC}; \cite[Section~3.37]{LangevinLean2026ULASC}.

\section{\texorpdfstring{Use the true covariance}{Use the true covariance}}


Apply the Gaussian covariance bilinear form to two linear observations. The standard Gaussian has mean zero and identity covariance, so the expectation of \ensuremath{\langle}\(u\),\(Z\)\ensuremath{\rangle}\ensuremath{\langle}\(v\),\(Z\)\ensuremath{\rangle} equals \ensuremath{\langle}\(u\),\(v\)\ensuremath{\rangle}. Square integrability justifies the product integral.


\section{\texorpdfstring{Turn a quadratic form into its trace}{Turn a quadratic form into its trace}}


Expand both arguments of a continuous bilinear form in the standard orthonormal basis. Interchange finite sums and integrals using Gaussian second moments. The covariance identity kills off-diagonal terms, leaving the sum of the diagonal values.


\section{\texorpdfstring{Average the entire Euler Taylor polynomial}{Average the entire Euler Taylor polynomial}}


Expand the increment \(hv+\sqrt{2h}Z\) in the linear and quadratic derivative terms. Gaussian centering removes its linear-noise and drift-noise terms; the trace is the actual Laplacian. The remaining deterministic quadratic drift is \(h^{2}D^{2}\varphi [v,v]/2\). This is an exact equality for the Taylor polynomial, prior to any remainder estimate.



Earlier results used here: \cite[Section~3.55.1]{LangevinLean2026ULASC}.

\begin{thebibliography}{ULA-LSI}
\bibitem[ULA-SC]{LangevinLean2026ULASC}
LangevinLean Project. \emph{Unadjusted Langevin under Strong Convexity}. LL-01, 2026. Draft report, content revision 263fb89cdbf0. \url{https://odakyildiz.github.io/LangevinLean/docs/pdf/UnadjustedLangevin/FullProof.pdf}.
\end{thebibliography}

\end{document}
